\topic{foundations}{Mathematical structure and physical assumptions}

These notes develop engineering models with the mathematical depth their claims
require. A successful computation is evidence about a specified model, initial
data and discretization; it is not a general existence theorem or experimental
validation of that model. Book sections are reading anchors, not claims that the
textbook states every extension developed here.

\subsection{Configuration, state and conventions}
A rigid body's configuration is a point of $\SO(3)$, or of $\operatorname{SE}(3)$
when translation is included. Its dynamical state additionally needs velocities
or momenta: configuration alone is not a complete six-degree-of-freedom state.
For $R:\R^3_{\rm body}\to\R^3_{\rm inertial}$, define
$\widehat\omega x=\omega\times x$. Then $\dot R=R\widehat\omega$ uses
\emph{body} angular velocity, whereas $\dot RR^T$ is its inertial representation.
Units, active/passive interpretation and multiplication order must accompany a
rotation API. Read Wie's inertia and angular-momentum development
\cite[Sections 6.1--6.4]{wie} alongside Ernest's geometric notes.

\subsection{From geometry to conservation}
On an oriented Riemannian manifold with metric $g$ and volume form $\mu_g$,
write $u^i$ for the velocity components in coordinates $x^i$ and define
divergence by
\[
 \mathcal L_u\mu_g=(\operatorname{div}_g u)\mu_g,
 \qquad
 \operatorname{div}_g u=\frac{1}{\sqrt{\det g}}
   \frac{\partial}{\partial x^i}\!\left(\sqrt{\det g}\,u^i\right).
\]
For smooth density $\rho$ and velocity $u$, conservation of mass on every
material region $V_t$ transported by $u$ gives
\[
 0=\frac{d}{dt}\int_{V_t}\rho\,\mu_g
   =\int_{V_t}\bigl(\frac{\partial\rho}{\partial t}+\operatorname{div}_g(\rho u)\bigr)\mu_g.
\]
Localization yields the continuity equation. In a thin duct, integrating over
cross sections and assuming stationary, effectively uniform section data gives
$\rho u A=\dot m$. Those reductions are physical assumptions, not consequences
of differential geometry alone. The nozzle study uses them explicitly
\cite[Sections 3.2--3.3]{sutton}.

\subsection{Regularity, shocks and honest evidence}
The pointwise argument presumes enough differentiability to apply transport
and localization. For shocks, formulate conservation weakly and impose an
appropriate entropy/admissibility condition. The smooth isentropic nozzle model
does not describe entropy production across a shock. Boundary data, constitutive
closure and any viscosity/heat-transfer approximation belong to the model.
Nothing here proves global smoothness of three-dimensional Navier--Stokes.

In each executable study distinguish: a mathematical identity and its domain;
a physical modeling assumption; a symbolic check; a numerical test at declared
resolution; and agreement with measurement. Numerical order and energy drift
are different properties \cite[Chapter 17]{nr}. Personal Reading Room checkboxes
record study progress independently of computation evidence.
